%%%%%%%%%%%%%%%%%%%%%%%%%%%%%%%%%%%%%%%%%%%%%%%%%%%%%%%%%%%%%%%%%%%%%%%%%%%%%%
% Introdução do capitulo %
%%%%%%%%%%%%%%%%%%%%%%%%%%%%%%%%%%%%%%%%%%%%%%%%%%%%%%%%%%%%%%%%%%%%%%%%%%%%%%
\label{cap:rsl}

A \acrfull{RSL} foi realizada de acordo com o protocolo definido por Kitchenham e Charters \cite{barbara2007guidelines}, com o objetivo de identificar o estado da arte sobre implementação de princípios éticos em sistemas de \acrshort{IA} aplicados à Engenharia de Software. A \acrshort{RSL} constitui processo estruturado de identificação, avaliação crítica e interpretação sistemática de evidências científicas relacionadas a uma área ou questão de pesquisa especifica \cite{Kitchenham2009}. Esse processo é dividido em três fases metodológicas distintas \cite{barbara2007guidelines}:

\begin{enumerate}
    \item \textit{Planejamento}: Definição da necessidade da revisão, identificação de objetivos específicos e elaboração do protocolo investigativo, compreendendo os seguintes artefatos: questões de pesquisa, \textit{string} de busca, critérios de seleção de estudos, lista de dados a serem extraídos e \textit{checklist} de avaliação de qualidade.
    
    \item \textit{Execução}: Realização do protocolo previamente estabelecido mediante aplicação sistemática dos artefatos desenvolvidos para triagem de estudos e extração estruturada de informações relevantes;
    
    \item \textit{Síntese}: Documentação analítica dos achados da revisão, materializada neste capítulo como componente fundamental da dissertação.
\end{enumerate}

%%%%%%%%%%%%%%%%%%%%%%%%%%%%%%%%%%%%%%%%%%%%%%%%%%%%%%%%%%%%%%%%%%%%%%%%%%%%%%

\section{Questões de Pesquisa}

O propósito desta \acrshort{RSL} é esclarecer o estado atual do conhecimento científico sobre práticas efetivas para incorporação de princípios éticos no desenvolvimento de sistemas computacionais baseados em \acrshort{IA}, com foco especial em abordagens operacionalizáveis no contexto da Engenharia de Software. 
O escopo da revisão se restringe apenas a contribuições que estabeleçam conexões diretas entre princípios éticos e práticas aplicáveis ao ciclo de vida de desenvolvimento, excluindo intencionalmente discussões puramente filosóficas, normativas ou que não estejam ancorados em processos de engenharia. 
A Tabela~\ref{tab:research_questions} apresenta as questões de pesquisa.

\begin{table*}[htb!]
  \centering
  \caption{Questões de pesquisa}
  \label{tab:research_questions}
  \begin{tabular}{cp{13cm}}
    \hline
    \textbf{Identificador} & \textbf{Questão de Pesquisa} \\
    \hline
    \textbf{RQ.1} & Quais princípios éticos específicos são identificados na literatura como aplicáveis ao desenvolvimento de sistemas de \acrshort{IA}?  \\
    \textbf{RQ.2} & Quais categorias de abordagens metodológicas (\textit{frameworks}, métodos, processos, ferramentas) são propostas para operacionalização de ética no ciclo de vida do software?  \\
    \textbf{RQ.3} & Quais evidências empíricas são reportadas quanto à efetividade, custos de implementação, limitações técnicas e riscos associados às abordagens de operacionalização? \\
    \textbf{RQ.4} & Quais lacunas de conhecimento, desafios não resolvidos e necessidades futuras de pesquisa são identificados na integração sistemática de ética ao \acrfull{SDLC}? \\
  \hline
\end{tabular}
\end{table*}

Para esta investigação foram considerados como sistemas baseados em \acrshort{IA} aqueles que empregam técnicas computacionais de aprendizado de máquina \textit{(machine learning)}, aprendizado profundo \textit{(deep learning)} ou implementam tomadas de decisão autônomas. 
A definição metodológica em torno de abordagens operacionalizáveis baseia-se no reconhecimento das diferenças significativas entre discursos éticos teórico-normativos e práticas concretas de implementação técnica, além da existência de estratégias distintas de aplicação de acordo com contextos específicos. 
Embora fundamentos éticos possam apresentar convergências conceituais entre diferentes domínios, é metodologicamente imprescindível examinar separadamente suas manifestações operacionais, pois não há garantia de que métodos de operacionalização sejam transferíveis entre contextos diferentes, o que expandiria substancialmente o escopo investigativo além dos limites viáveis para este estudo.

%%%%%%%%%%%%%%%%%%%%%%%%%%%%%%%%%%%%%%%%%%%%%%%%%%%%%%%%%%%%%%%%%%%%%%%%%%%%%%

\section{String de Busca}

A construção da \textit{string} de busca baseou-se no método \textbf{PICOC} (População, Intervenção, Comparação, Resultado, Contexto) proposto por Wohlin \cite{Wohlin2012}. Neste método estruturado, \textbf{População} designa o objeto primário de estudo; \textbf{Intervenção} corresponde aos meios ou mecanismos aplicados para alcançar determinado objetivo; \textbf{Comparação} identifica elementos confrontados com a intervenção (quando aplicável); \textbf{Resultado} especifica as consequências observáveis da intervenção; e \textbf{Contexto} delimita o foco investigativo, incluindo suas restrições e condições de contorno. A Tabela~\ref{tab:picoc} apresenta a especificação dos termos PICOC empregados na formulação da string de busca genérica. O elemento Comparação não se aplica ao presente estudo, pois não há confrontação entre diferentes intervenções.

\begin{table}[htbp]
  \centering
  \caption{Especificação dos termos PICOC}
  \label{tab:picoc}
  \begin{tabular}{llp{9cm}}
    \hline
    \textbf{PICOC} & \textbf{Termo Principal} & \textbf{Termos Relacionados e Sinônimos} \\
    \hline
    População & ética &	inteligência artificial ética, IA responsável, IA confiável, accountability, transparência algorítmica, justiça algorítmica, viés, privacidade, explicabilidade   \\
    Intervenção  & inteligência artificial 	& IA, aprendizado de máquina, machine learning, ML, aprendizado profundo, deep learning, sistemas autônomos \\
    Comparação & \textit{não aplicável} & \textit{não aplicável} \\
    Resultado & engenharia de software 	& desenvolvimento de software, engenharia de requisitos, DevOps, MLOps, ciclo de vida de software, SDLC, desenvolvimento ágil \\
    Contexto & operacionalização & framework, método, processo, ferramenta, diretriz, abordagem, padrão, modelo  \\
  \hline
\end{tabular}
\end{table}

O processo de refinamento da estratégia de busca iniciou-se com a identificação de termos-chave principais para cada dimensão PICOC: \textit{ética}, \textit{inteligência artificial}, \textit{engenharia de software} e \textit{operacionalização}. Subsequentemente, realizou-se investigação exploratória iterativa com termos correlatos e sinônimos técnicos para validação da abrangência e relevância dos resultados recuperados. Após múltiplas iterações de refinamento e validação com estudos de controle conhecidos \cite{Agbese2023,Vakkuri2022,Cerqueira2021}, estabeleceu-se a string de busca genérica final:

\vspace{0.3cm}
\colorbox{formalshade}{\parbox{0.9\columnwidth}{\small\emph{(ethic* OR ``ethical AI'' OR ``responsible AI'' OR ``trustworthy AI'' OR accountability OR transparency OR fairness OR bias OR privacy OR explainab* OR interpretab*) \textbf{AND} (``artificial intelligence'' OR AI OR ``machine learning'' OR ML OR ``deep learning'') \textbf{AND} (``software engineering'' OR ``software development'' OR ``requirements engineering'' OR DevOps OR MLOps OR ``software lifecycle'' OR SDLC) \textbf{AND} (framework OR method* OR process OR tool OR guideline OR governance OR model OR approach)}}}
\vspace{0.3cm}

As bases bibliográficas digitais selecionadas para execução da estratégia de busca compreendem \href{https://dl.acm.org/}{ACM Digital Library}, \href{https://ieeexplore.ieee.org/}{IEEE Xplore}, \href{https://www.scopus.com/}{Scopus} e \href{https://www.webofscience.com/}{Web of Science}. A seleção fundamentou-se em três critérios objetivos: (i) relevância consolidada para pesquisas em Engenharia de Software \cite{Brereton2007}; (ii) abrangência de indexação de conferências e periódicos de alto impacto na área; e (iii) capacidade técnica dos mecanismos de busca para processamento integral da string genérica formulada. A Tabela~\ref{tab:strings} especifica as adaptações sintáticas da string genérica para cada fonte bibliográfica, respeitando particularidades dos respectivos mecanismos de busca.

\begin{table}[htbp]
 \centering
 \caption{\textit{strings} de Busca por Fonte}
 \label{tab:strings}
 \footnotesize
 \begin{tabular}{p{2cm}p{11.5cm}}
    \hline
    Fonte & \textit{string} \\
    \hline
    \addlinespace
    \textbf{ACM Digital Library} & Abstract: ("ethic*" OR "ethical AI" OR "responsible AI" OR "trustworthy AI" OR accountability OR transparency OR fairness OR bias OR privacy OR explainab* OR interpretab* OR "algorithmic accountability") Abstract: ("artificial intelligence" OR AI OR "machine learning" OR ML OR "deep learning") Abstract: ("software engineering" OR "software development" OR "requirements engineering" OR DevOps OR MLOps OR "software lifecycle" OR SDLC OR "development lifecycle" OR "user stories") Abstract: (framework OR method* OR process OR tool OR checklist OR pattern OR guideline OR governance OR model OR approach)  \\
    \hline
    \addlinespace
    \textbf{IEEE Xplore} & 
    ("Abstract":"ethic*" OR "ethical AI" OR "responsible AI" OR "trustworthy AI" OR accountability OR transparency OR fairness OR bias OR privacy OR explainab* OR interpretab* OR "algorithmic accountability") AND ("Abstract":"artificial intelligence" OR AI OR "machine learning" OR ML OR "deep learning") AND ("Abstract":"software engineering" OR "software development" OR "requirements engineering" OR DevOps OR MLOps OR "software lifecycle" OR SDLC OR "development lifecycle" OR "user stories") AND ("Abstract":framework OR method* OR process OR tool OR checklist OR pattern OR guideline OR governance OR model OR approach) \\
    \hline
    \addlinespace
    \textbf{Scopus} & 
    TITLE-ABS-KEY (("ethic*" OR "responsible AI" OR "trustworthy AI" OR accountability OR transparency OR fairness OR bias OR privacy OR explainab* OR interpretab* OR "algorithmic accountability") AND ("artificial intelligence" OR AI OR "machine learning" OR ML OR "deep learning") AND ("software engineering" OR "software development" OR "requirements engineering" OR DevOps OR MLOps OR "software lifecycle" OR SDLC OR "development lifecycle" OR "user stories") AND (framework OR method* OR process OR tool OR checklist OR pattern OR guideline OR governance OR model OR approach)) AND PUBYEAR > 2019 AND PUBYEAR < 2026 AND PUBYEAR > 2020 AND PUBYEAR < 2026 AND (LIMIT-TO (PUBSTAGE,"final")) AND (LIMIT-TO (DOCTYPE,"cp") OR LIMIT-TO (DOCTYPE,"ar")) AND (LIMIT-TO (SUBJAREA,"COMP") OR LIMIT-TO (SUBJAREA,"ENGI") OR LIMIT-TO (SUBJAREA,"DECI") OR LIMIT-TO (SUBJAREA,"SOCI")) AND (LIMIT-TO (LANGUAGE,"English")) \\
    \hline
    \addlinespace
    \textbf{Web of Science} & 
    TS=(("ethic*" OR "ethical AI" OR "responsible AI" OR "trustworthy AI" OR accountability OR transparency OR fairness OR bias OR privacy OR explainab* OR interpretab* OR "algorithmic accountability") AND ("artificial intelligence" OR AI OR "machine learning" OR ML OR "deep learning") AND ("software engineering" OR "software development" OR "requirements engineering" OR DevOps OR MLOps OR "software lifecycle" OR SDLC OR "development lifecycle" OR "user stories") AND (framework OR method* OR process OR tool OR checklist OR pattern OR guideline OR governance OR model OR approach)) AND PY=(2020-2025) \\
    \bottomrule
    \end{tabular}
    \end{table}

A ideia inicial era executar as \textit{strings} para título, resumo e palavras-chave. No entanto, os mecanismos de busca não compartilhavam essa opção específica. Na ACM, era ou o título ou o resumo sem repetir a \textit{string}. Por isso escolhemos o resumo. No IEEE, \textit{Abstract} refere-se a título, resumo e palavras-chave, assim como \textit{TITLE-ABS-KEY} no Scopus. No Web of Science, \textit{TS} (Topic Search) busca em título, resumo, palavras-chave do autor e Keywords Plus.

Foram utilizadas as opções de filtragem e \textit{string} disponíveis para cada fonte para aplicar alguns de nossos critérios de exclusão e acelerar a seleção de artigos. Por exemplo, na ACM Library, aplicamos filtro manual para \textit{Content Type} selecionando apenas \textit{Research Article} e selecionando um intervalo de datas de 2020 até o momento da busca. No IEEE Xplore, aplicamos filtro manual para o intervalo de datas. No Scopus, todos os filtros foram adicionados à \textit{string}: o \textit{Document Type} como \textit{Conference Paper} ou \textit{Article}; o \textit{Publication Year} como posterior a 2019; \textit{Language} como \textit{English}; e o \textit{Source Type} como \textit{Journal} ou \textit{Conference Proceedings}. No Web of Science, o filtro de ano foi aplicado diretamente na \textit{string}.

%%%

